\chapter*{Conferencia XI}
\addcontentsline{toc}{chapter}{LECTURA I - El timbre y el carácter del violín}
SEÑORES: Se reanuda la disminución del volumen y la intensidad del sonido del violín en el puente.
(b) El grosor del puente puede ser tan grande que se convierte en un factor importante para disminuir tanto el volumen como la intensidad del sonido. Demostrar este hecho es muy sencillo. Si bien este hecho es algo que los estudiantes de violín observan a diario, dado que solo una pequeña minoría de ellos se preocupa por la filosofía relacionada con los modificadores de tono, y puesto que el conocimiento de ciertas leyes físicas es valioso para el regulador de tono del violín, parece aconsejable considerar dichas leyes. Ciertas leyes físicas son valiosas para el sonido del violín en la medida en que el ingenio humano puede aplicarlas. En dicha aplicación reside una dificultad. En toda mi larga experiencia, solo he obtenido un resultado de valor significativo para el sonido del violín siguiendo un razonamiento a priori.
Considero que el problema del grosor del puente es complejo y tiene cuatro factores inevitables:
(a) Altura del diapasón. (b) Altura del arco. (c) Diámetro de las cuerdas. (d) Densidad de la fibra.
(a) Como se mostró anteriormente, la altura del puente está determinada por la altura del diapasón; sin embargo, esta afirmación está sujeta a modificación. El hecho de que dos violines tengan una altura de diapasón exactamente similar no implica necesariamente que los puentes tengan una altura exactamente similar. Este hecho se debe a la disimilitud en la altura del arco.
171

en la posición del puente. Sin duda, muchos, si no todos los estudiantes de violín, han observado que algunos fabricantes colocan el punto más alto del arco en la posición del puente, mientras que otros lo colocan más arriba o hacia adelante desde la posición del puente. Es evidente que esta diferencia requiere una diferencia en la altura del puente. [Aunque no viene al caso, el temor a omitir algo me obliga a llamar la atención sobre la colocación del punto más alto del arco (arco longitudinal) en la posición del puente. Según mi observación, colocar el punto más alto del arco en la posición del puente aumenta la potencia del sonido. Desde mi punto de vista, las razones de dicho aumento son dos: Primero: la disminución de la altura del puente permite la disminución del grosor del puente; por lo tanto, la disminución de los efectos de amortiguación debido a una mayor masa de madera en el puente. Segundo:
Colocar el punto más alto de la placa arqueándose hacia delante o hacia arriba de la posición del puente reduce la amplitud de la oscilación de la caja de resonancia debajo de las cuerdas; por lo tanto, disminuye la potencia del tono.
Es evidente que el grosor del puente puede amortiguar el sonido del violín. También es evidente que la disminución de la altura del puente permite la disminución de su grosor. Se observa que la disminución del grosor del puente, hasta el punto de mantener la presión de las cuerdas sin doblarlas, ejerce sobre la caja de resonancia una mayor fuerza de los golpes de las cuerdas. Por lo tanto, tanto la disminución de la altura como la disminución del grosor del puente pueden amortiguar el sonido del violín.
172

Erar para disminuir la potencia del tono.
[En una hora anterior se demostró que la proximidad de cuerpos vibrantes, susceptibles a una fuerza idéntica, opera para aumentar la acción simpática. Por esta razón, la pérdida de acción simpática entre las cuerdas y esa parte de la caja de resonancia debajo de las cuerdas podría considerarse como uno de los modificadores que operan para disminuir el volumen y la intensidad del tono del violín; pero debido a que la lista de tales modificadores de tono ya llega al número 12, y, debido a que no se encontraron dos modificadores más de este tipo, por lo tanto, me niego a permitir que la acción simpática, (con su número 13) tenga la oportunidad de "hacer magia negra" mi
trabajar.
(c) El diámetro de las cuerdas del violín determina en gran medida (aunque no totalmente) la presión descendente sobre el puente. Es evidente que un diámetro de cuerda menor produce menos presión descendente sobre el puente debido a la menor tensión requerida para la afinación. Por lo tanto, para producir la máxima potencia tonal con cuerdas más delgadas se requiere una disminución del grosor del puente proporcional a la disminución del diámetro de la cuerda. La excepción a la regla del diámetro de la cuerda sobre la presión descendente sobre el puente reside en el arqueamiento de la placa. Es evidente que, sin arqueamiento alguno, la presión descendente de las cuerdas es mínima; y, debido a que las cuerdas se encuentran a la menor distancia práctica por encima de una línea recta desde la selleta hasta la cejuela. En dicha línea recta, siendo la presión descendente cero, se deduce que cada grado de altura del puente por encima de esta línea aumenta la presión descendente de las cuerdas.

[*]Por lo tanto, cuanto mayor sea el arqueo de la placa, mayor será la presión descendente de las cuerdas sobre el puente. Por consiguiente, el arqueo se convierte en un factor determinante del grosor del puente. Así se demuestra que no se pueden aplicar reglas estrictas al grosor del puente del violín. (c) La distancia entre los pedestales o vástagos del puente, si bien no es un factor determinante en la disminución de la potencia tonal, sí es un factor que requiere atención. Al ser demasiado grande o demasiado pequeña, esta distancia reduce la amplitud de la oscilación de la caja de resonancia y, por lo tanto, disminuye la potencia tonal.
Hay dos factores que rigen el tramo del puente: la posición de la barra.
Distancia entre salidas.
Según mi observación, no colocar el centro del pedestal izquierdo sobre el centro del grosor de la barra reduce la potencia tonal de la cuerda D. De nuevo, colocar el centro del pedestal derecho sobre las fibras de la caja de resonancia cortadas por la salida derecha reduce la potencia tonal de la cuerda E. El puente, sin pedestales, tiene interés, por lo tanto: en aquellos casos en los que un tono hueco de gran volumen, pero de intensidad debilitada, causado por una reducción excesiva de la rigidez de la caja de resonancia en su área central, el puente sin pedestales, cuidadosamente ajustado a la curva lateral de
La caja de resonancia funciona para aumentar la potencia tonal de las cuerdas A y D; no mucho, pero sí perceptiblemente. Sobre esa caja de resonancia que posee suficiente
rigidez en su área central, no encuentro que la
puente sin pedestales funciona hacia ago- 174

disminución o reducción de la potencia tonal de cualquier cuerda. En este sentido, encuentro que el puente sin pedestales funciona exactamente como el bloque de refuerzo pegado transversalmente sobre la superficie interior de la tapa armónica en la posición del puente.
[En el trabajo experimental, he colocado dicho bloque de refuerzo sobre 17 diferentes cajas de resonancia con distintos grados de rigidez en sus áreas centrales. El beneficio para el tono de ello solo se manifestó sobre aquellas cajas de resonancia cuyo grosor se redujo lo suficiente como para causar debilitamiento.
Tonos A y D. En aquellas cajas de resonancia que producían tonos A y D fuertes, el bloque de refuerzo no disminuía ni aumentaba la potencia sonora de ninguna cuerda.
La densidad de la fibra es un factor poderoso en el puente. Tanto la fibra más blanda como la más densa tienden a disminuir la potencia del sonido. Sin duda, este hecho se debe a una falla en la transmisión de la fuerza de las cuerdas a la caja de resonancia. En el caso del violín, es evidente que la fuerza de las cuerdas solo puede comunicarse a la caja de resonancia mediante la acción vibratoria de los medios de conexión, como el puente y el aire. Dicha transmisión no se limita al puente de ninguna manera. La acción simpática entre las cuerdas y la caja de resonancia se debe completamente a la presencia del aire como medio de conexión. Considero que la potencia de la acción simpática es sumamente digna de consideración por parte del estudiante de violín. No conozco una demostración más sencilla de dicha potencia que la siguiente: Retire el puente y, sobre el bloque de la cola, coloque un puente de altura suficiente para
175

Mantén las cuerdas a su altura habitual sobre el diapasón.
[Para realizar una prueba precisa de la potencia en la acción simpática, es una condición necesaria que las cuerdas no estén unidas al violín. En ambos métodos, la potencia tonal de la acción simpática es sorprendentemente grande cuando la rigidez de la caja de resonancia corresponde a los diámetros de las cuerdas, en otras palabras, a la fuerza de las cuerdas; pero, cuando dicha rigidez es demasiado grande, entonces la potencia del tono simpático disminuye, y disminuye porque la acción simpática máxima entre dos cuerpos vibrantes contiguos exige una susceptibilidad igual a la fuerza. Dicho tono simpático también disminuye por el área de las salidas.] La mayor densidad en cualquier puente que he probado se encuentra en uno hecho de hueso. La densidad en este puente modifica el tono de una manera peculiar; tanto el volumen como la intensidad del tono disminuyen enormemente, mientras que el poco tono restante es notablemente
delgada> y de una cualidad dolorosamente punzante.
El puente de fibra más suave que he probado está hecho de pino blando de Michigan. Este puente reduce tanto el volumen como la intensidad del tono, pero también reduce la calidad desagradable del mismo. En violines con un tono ruidoso, no he observado ninguna mejora en la calidad desagradable del tono tras el uso de esta madera blanda para el puente. El grosor de dicho puente no se puede reducir al mismo grado que el puente de arce debido a su mayor facilidad de flexión bajo la presión de las cuerdas. La aspereza del tono se puede reducir considerablemente empleando este tipo de madera blanda.
176

Madera tanto para el puente como para el poste.
El trabajo de volutas en el puente tiene una misión más importante que la mera ornamentación. En mi experiencia, el puente "completo" o "sólido", de cualquier madera, contribuye a disminuir la potencia tonal. Si se pudiera disminuir el grosor del puente "sólido" hasta que su acción vibratoria fuera igualmente susceptible a la fuerza que la del puente tallado en volutas, entonces el trabajo de volutas se convertiría en meramente ornamental. Pero, tal disminución del grosor del puente es impracticable porque la presión descendente de las cuerdas tiende a doblar el puente "sólido" cuando se reduce su grosor. Es un hecho que doblar el puente disminuye la amplitud de su oscilación; por lo tanto, doblar el puente disminuye la potencia tonal, y, exactamente de la misma manera que doblar la caja de resonancia. Obviamente, el grosor del puente debe ser lo suficientemente grande como para mantenerlo erguido bajo la presión de las cuerdas. Debido a que la presión de las cuerdas varía con la altura del arco de la placa y los diámetros de las cuerdas, y debido a que diferentes muestras de madera para puentes presentan diferentes grados de rigidez, por lo tanto, reglas estrictas para
El grosor del puente no se aplica.	En este asunto, yo
No conozco ninguna regla exitosa aparte de la habitual.	
regla,	''Corta y prueba.''
"
La experiencia lo demuestra
que el puente "sólido", lo suficientemente grueso como para mantenerse erguido bajo una presión de cuerda de 25 a 28 libras, posee masa suficiente para amortiguar apreciablemente el tono. La experiencia también demuestra que parte de dicha masa puede eliminarse de forma segura de ciertas partes del puente sin disminuir la rigidez.
177

hasta el punto de doblarse.
Dicha disminución de masa se consigue mediante el familiar trabajo de volutas:
La ubicación y extensión de los adornos en espiral del puente no deben dejarse al azar, ya que una mala posición o extensión de estos adornos puede provocar irregularidades en la potencia del sonido. Por lo tanto, incluso después de haber prestado la máxima atención a cada detalle de la construcción que contribuye a la uniformidad del sonido, la mala posición y extensión de los adornos centrales en espiral del puente pueden suponer un reto incluso para el constructor de violines más hábil.
Mientras sostengo este puente frente a usted, observa la ubicación y extensión de la voluta central, y también el trabajo de voluta en cada extremo. Cerca de cada extremo, observa un istmo, o cuello estrecho entre las volutas central y de los extremos. Como puede observar, estos cuellos conectan las mitades superior e inferior del puente. Es evidente que toda la acción vibratoria en el puente, provocada por la acción de las cuerdas, debe viajar hacia abajo a través de estos cuellos estrechos. Es evidente que la disminución de estos cuellos modifica la transmisión de la acción vibratoria. También es evidente que la desigualdad en las dimensiones de estos cuellos produce una susceptibilidad desigual a la fuerza. También es evidente que la fuerza en las cuerdas del violín varía según su diámetro y peso. Por lo tanto, resulta obvio que la igualdad en las dimensiones de estos cuellos disminuye la potencia tonal de las cuerdas más pequeñas. Por consiguiente, para producir una potencia tonal uniforme, las dimensiones...
178

las secciones del mástil debajo de las cuerdas A y E
debe ser menor que las dimensiones del mástil debajo de las cuerdas G y D; y dicha disminución de masa debe estar en la misma proporción que los diámetros disminuidos de las cuerdas.
Aquí hay una docena de puentes de alta calidad. Con una regla de maquinista, determino las dimensiones de sus istmos entre el trabajo de voluta central y el extremo. Encuentro que el grosor de esos istmos es de 1 pulgada; pero, en su ancho, encuentro una variación, con respecto a la igualdad, de 1 a 16 pulgadas. Es evidente que dicha variación en la masa puede aprovecharse para asegurar una uniformidad en la potencia del sonido. Así, al colocar el istmo más grande debajo de las cuerdas G y D, se logra igualar la transmisión de fuerza. También es evidente que la disminución del grosor de estos istmos permite que el puente se doble bajo la presión de las cuerdas; pero, la disminución de su ancho no puede causar un resultado tan desastroso. También es evidente que las dimensiones de esos istmos deben regirse por la presión de las cuerdas; y, dado que dicha presión varía con diferentes grados de curvatura de la placa y diámetro de la cuerda, por lo tanto, para tales dimensiones, no se pueden aplicar reglas estrictas.
(f) La madurez de la madera de puente (madurez del árbol antes de ser talado) es un factor importante en la lista de modificadores del tono. Dicha madurez es tan importante como la madurez de la madera de la caja de resonancia. La ausencia de madurez en cualquiera de ellas conlleva una pérdida de la calidad del sonido. Las razones para tomar la madera de la caja de resonancia de la parte del tronco comprendida entre el duramen y la albura se aplican igualmente a la madera de puente. 179
